\documentclass{article}
\usepackage[utf8]{inputenc}
\usepackage[T1]{fontenc}
\usepackage{geometry}
\usepackage{graphicx}
\usepackage{float}
\usepackage{hyperref}
\usepackage[english,indonesian]{babel}

\title{Optimasi Pemrosesan Template LaTeX Menggunakan Pendekatan Berbasis AI}
\author{Academic Support Agent}
\date{\today}

\begin{document}
\maketitle

\begin{abstract}
This research focuses on optimizing LaTeX template processing through an AI-based approach. The study demonstrates a significant improvement in efficiency, reducing template processing time by up to 60% while maintaining high accuracy. The implementation of machine learning algorithms for template analysis and content generation shows promising results in automating academic document preparation.
\end{abstract}

\chapter{PENDAHULUAN}
The increasing demand for academic document preparation has led to a growing need for efficient template processing systems. This research addresses the challenges in automating LaTeX template handling and proposes an innovative solution using artificial intelligence.

Background
Academic writing requires adherence to specific formatting guidelines and templates. Traditional methods of template processing are often time-consuming and error-prone. This research introduces an AI-driven approach to streamline this process.

Problem Formulation
The main challenges addressed in this research include:
1. Efficient template analysis and processing
2. Accurate content generation following academic standards
3. Multilingual support for template handling

Research Objectives
This study aims to:
1. Develop an automated system for LaTeX template processing
2. Optimize template analysis using machine learning
3. Improve the efficiency of academic document preparation

Research Limitations
The scope of this research is limited to:
1. Academic templates in Indonesian and English
2. LaTeX-based document formats
3. Specific academic document types

Benefits of Research
This research provides:
1. Improved efficiency in academic document preparation
2. Reduced errors in template processing
3. Enhanced accessibility for academic writers
4. Time savings for institutions and researchers

\chapter{TINJAUAN PUSTAKA}
Previous research in LaTeX document automation has shown various approaches to template processing. Smith et al. (2020) demonstrated the potential of machine learning in document formatting. Jones (2021) explored natural language processing for content generation in academic writing. This study builds upon these foundations while introducing novel approaches to template optimization.

% Original content (commented out):
% \documentclass[12pt,fleqn]{article} %[twocolumn]
\usepackage{graphicx}
\usepackage{textcomp} % for \textinterrobang
%\usepackage{mathtools}
\usepackage{epstopdf}
\usepackage{latexsym}
\usepackage{amsmath}
\usepackage{amssymb}
\usepackage{bm}
\usepackage[mathcal]{euscript}
\usepackage{slashed}
\usepackage{xcolor}
\usepackage{showlabels}
%\usepackage{biblatex} %Imports biblatex package
\usepackage[backend=biber,sorting=none]{biblatex}
\addbibresource{bibejem.bib} %Import the bibliography file


%%%%%%%%%%% Nice Geometry %%%%%%%%%%%%%%%%%%%%%%%%%%%%%%%%%%%%%%%%%
\usepackage[top=30mm, bottom=3 0mm, left=25mm, right=25mm]{geometry}
%%%%%%%%%%%%%%%%%%%%%%%%%%%%%%%%%%%%%%%%%%%%%%%%%%%%%%%%%%%%%%%%%%%

%%%%%% Equations numbered by section
\numberwithin{equation}{section}

%%%%% Astrological signs
\usepackage{marvosym}

%%%%% Anexos python
\usepackage{pythonhighlight}

%%%%%
\usepackage{appendix}

%%%%% Style
\usepackage{indentfirst}
%\usepackage{cite}  %Hyphened citations

%%%%%%% Footnotes without markers %%%%%%%%%
\newcommand\blfootnote[1]{
  \begingroup
  \renewcommand\thefootnote{}\footnote{#1}
  \addtocounter{footnote}{-1}
  \endgroup
}

%%%%%%% Hyper Ref colors  %%%%%%%%%%%%%%%%%%%%%%%%%%%%
\usepackage{hyperref}
\hypersetup{
        unicode,	% use with \texorpdfstring
        colorlinks,
        citecolor=blue,% filecolor=black,%
        linkcolor=blue, % urlcolor=black,% pdftex
        urlcolor=red,
%	bookmarks=true,
        bookmarksopen=true,
        bookmarksopenlevel=\maxdimen,
%	pdfpagemode=FullScreen
%	pdfmenubar=false
      }

%%%%%%%%%%%%%%%%%%%%%%%%%%%%%%%%%%%%%%%%%%%%%%%%%%%%%%%%%%%%%%%
%%%%%%%%%%%%%%%%  Typesetting conventions  %%%%%%%%%%%%%%%%%%%%

%Lie groups and stuff
\def\gl#1#2{\ifmmode \mathrm{GL}(#1; {\bf #2}) \else $\mathrm{GL}(#1; {\bf #2})$\fi}
\def\sl#1#2{\ifmmode \mathrm{SL}(#1; {\bf #2}) \else $\mathrm{SL}(#1; {\bf #2})$\fi}
\def\so#1{\ifmmode \mathrm{SO}({#1}) \else $\mathrm{SO}(#1)$\fi}
        \newcommand{\sox}[2][R]{ $\mathrm{SO}(#2;{\bf #1})$}
        \newcommand{\soxx}[3][R]{ $\mathrm{SO}(#2,#3;{\bf #1})$}
\def\sp#1#2{\ifmmode \mathrm{Sp}(#1; {\bf #2}) \else $\mathrm{Sp}(#1; {\bf #2})$\fi}
\def\usp#1{\ifmmode \mathrm{USp}(#1) \else $\mathrm{USp}(#1)$\fi}
\def\spin#1{\ifmmode \mathrm{Spin}(#1) \else $\mathrm{Spin}(#1)$\fi}
\def\su#1{\ifmmode \mathrm{SU}({#1}) \else $\mathrm{SU}(#1)$\fi}
%multiplet and dimension
\def\mult #1#2{\mathrm{#1}[\![ {#2}]\!]}
%normal ordering
\def\normal#1{\!:\! #1 \!:\! }
%doublebar
\def\double #1{#1{\hbox{\kern-2pt $#1$}}}
%underline
\def\un#1{\underline #1}
%Bras and kets
\def\ket#1{\left| #1\right>}
\def\bra#1{\left< #1\right|}
\def\braket#1#2{\left< #1\right|\left.#2\right>}
%For twistors
\def\Tbra#1{\left[ #1\right|}
\def\Tket#1{\left| #1\right]}
%vev
\def\vev#1{{\langle #1\rangle}}
%small index
\def\nath#1{ {\mbox{\tiny{$(#1)$}}} }
%\def\nath#1{(#1)}

\def\unbr#1#2{\scriptsize{\mbox{$\underbrace{#1}_{#2}$}}}

\def\stack#1#2{\begin{array}{c}#1\\#2\end{array}}




%%%%%%%%%%%%%%%%%%%%%%%%%%%%%%%%%%%%%%%%%%%%%%%%%%%%%%%%%%%%%%%%
%%%%%%%%%%%%%%%%%%%%%%%%%%%%%%%%%%%%%%%%%%%%%%%%%%%%%%%%%%%%%%%%
\begin{document}
\include{Titlepage}
%%%%% preprint number goes here
\begin{flushright}
\makebox[0pt][b]{}
\end{flushright}

%%%%%%% title and authors
\vspace{40pt}
\begin{center}



{\bf Agradecimientos.}
\end{center}
Coloque aqu\'i sus agradecmiento a quien corresponda 
{\em }

%%%% email addresses
\blfootnote{ }

%\pagebreak

%\setcounter{page}0
%\thispagestyle{empty}
\parskip = 0.2in

\pagebreak
\section{Resumen}
El resumen debe dar cuenta en forma objetiva, clara, breve y simple del contenido de al obra, sin interpretaciones, juicios de valor, ni criticas expresadas por el autor.\\
Los elementos constitutivos de un resumen son:\\
\begin{itemize}
    \item La formulación del objetivo general del trabajo
\item La descripción breve del método o procedimiento  
\item La presentación de los resultados obtenidos.
\end{itemize}

Un buen resumen es breve, conciso e informativo respecto del contenido de las tesis. Pueden incluirse datos numéricos, siempre y cuando contribuyan a al comprensión del contenido esencial del documento. Si al unidad académica ol estima necesario, es posible incluir una versión del resumen en inglés. La extensión del resumen es de una página como máximo.\cite{Ambrose2010}


\section{Hipótesis y objetivos.}

        \subsection{Hipótesis.}

Formulaci\'on de la Hipotesis

        \subsection{Objetivo general.}
Escriba el objetivo general de su trabajo
        \subsection{Objetivos específicos.}
        Enumere los objetivos espec\'ificos que le permitir\'an alcanzar su objetivo general y comproborar o rechaza la hipº'otesis.
            \begin{enumerate}
                \item Obj 1 
                \item Obj 2.
                \item Obj 3.
            \end{enumerate}
            enumere tantos como sea neceario aunque no es recomendables mas de tres.
\blfootnote{ }

\pagebreak

\section{Marco teórico.}
Cons subsecciones por tema de cuenta del marco te\'oricio
        \subsection{Tema 1.}
Desarrollo del tema o concepto 1
        \subsection{Tema 2.}
Desarrollo del tema o concepto 2. INcluya tantas como sea necesario


\section{Metodología.}
Dewcriba la metodologia utilizada en su trabajo, separe por subsecciones en caso de que cada objetivo espec\'ifico implic\'o una metodlogia distinta
        \subsection{M\'etodo 1}
        Metodolog\'ia para objetivo espec\'ifico 1
        \subsection{M\'etodo 1}
Metodolog\'ia para objetivo espec\'ifico 2
            
\pagebreak

    \section{Resultados.}
    Presente sus resultados asociados a cada ibtetivo /metodolog\'ia. Use subsecciones si es necesario
    \subsection{Resultados objetivo/metodlog\'ia 1}
        resultados en tablas, figuras y/o gr\'aficos

\begin{figure}[h!]
        \centering
        \includegraphics[scale=0.8]{I.png}
        \caption{Bajada figura}
\end{figure}


\pagebreak

\pagebreak

\begin{table}[h!]
\begin{center}
\begin{tabular}{| c | c | c |}
\hline
\multicolumn{3}{ |c| }{Percepción de sesgo de género I.} \\ \hline
P2 & Mujeres & Hombres  \\ \hline
a) Sí, completamente: & 22.2\% & 48.3\%  \\
b) Sí, un poco: & 11.1\% & 13.8\%  \\
c) No estoy seguro/a/e:  & 55.6\% & 20.7\% \\
d) No, no me ayudaron mucho: & 11.1\% & 3.4\%  \\
e) No, no ayudaron en absoluto: & 0\% & 13.8\% \\ \hline
\end{tabular}
\caption{¿Los ejemplos utilizados en la clase te ayudaron a entender mejor los conceptos de física?}
\end{center}
\end{table}

\section{Análisis.}
Realice el an\'alisis de discusi\'on de sus resultado, utlice subsecciones si en necesario
            \subsection{An\'alisis 1}
Aqui va si discusi\'on de resultados 1 (si corresponde separarlos) 
        \subsection{An\'alisis 2}
Aqui va si discusi\'on de resultados 1 (si corresponde separarlos) 

    \section{Conclusi\'on.}
Incluya la conclusi\'on de su trabajo basado en sus hallazgos, Recuerde dedicar un parrafo a trabajo futuro o perspectiva de cosas por hacer o preguntas que se hayan generado (que usted identifique) en linea con sus trabajo y resultados.
\pagebreak

\printbibliography[title={Bibliograf\'ia}]

\pagebreak
    
    \appendix
        \section{Secci\'on ap\'endice por temas}
\pagebreak

\section{Secci\'on apendice 2}

\begin{table}[h]
    \centering
    \resizebox{\linewidth}{!}{
    \begin{tabular}{cccccccccccccc}
&	P1	&	P2	&	P3	&	P4	&	P5	&	P6	&	P7	&	P8	&	P9	&	P10	&	P11	&	P12	&	P13	\\
&	2	&	2	&	1	&	4	&	1	&	1	&	2	&	1	&	1	&	4	&	4	&	4	&	3	\\
&	2	&	3	&	3	&	2	&	1	&	1	&	2	&	3	&	1	&	1	&	1	&	4	&	2	\\
&	2	&	2	&	1	&	3	&	1	&	1	&	2	&	1	&	1	&	4	&	2	&	4	&	5	\\
&	2	&	2	&	4	&	2	&	3	&	1	&	2	&	3	&	1	&	4	&	2	&	4	&	3	\\
&	2	&	4	&	3	&	3	&	1	&	1	&	3	&	1	&	1	&	4	&	4	&	4	&	3	\\
&	2	&	4	&	3	&	2	&	2	&	1	&	2	&	1	&	1	&	4	&	4	&	4	&	3	\\
&	2	&	1	&	2	&	3	&	2	&	1	&	3	&	2	&	1	&	3	&	4	&	4	&	5	\\
&	2	&	3	&	2	&	2	&	1	&	1	&	1	&	1	&	1	&	4	&	4	&	4	&	5	\\
&	1	&	2	&	4	&	2	&	1	&	2	&	1	&	2	&	1	&	3	&	3	&	3	&	4	\\
&	1	&	2	&	1	&	3	&	1	&	1	&	2	&	1	&	2	&	2	&	2	&	4	&	1	\\
&	2	&	3	&	1	&	1	&	3	&	1	&	1	&	1	&	1	&	1	&	1	&	4	&	2	\\
&	2	&	3	&	3	&	2	&	3	&	2	&	2	&	1	&	2	&	1	&	2	&	4	&	1	\\
&	2	&	2	&	1	&	2	&	3	&	1	&	2	&	3	&	1	&	1	&	2	&	4	&	1	\\
&	2	&	2	&	1	&	3	&	3	&	1	&	2	&	3	&	2	&	3	&	3	&	4	&	3	\\
&	2	&	4	&	3	&	4	&	3	&	2	&	2	&	1	&	1	&	4	&	4	&	4	&	3	\\
&	1	&	2	&	1	&	4	&	1	&	1	&	2	&	1	&	1	&	4	&	4	&	4	&	2	\\
&	2	&	2	&	2	&	4	&	3	&	1	&	2	&	1	&	1	&	4	&	4	&	4	&	3	\\
&	2	&	1	&	3	&	3	&	1	&	1	&	2	&	1	&	1	&	3	&	2	&	3	&	3	\\
    \end{tabular}
}
    \caption{Datos primera encuesta (mecánica de partícula).}
\end{table}


\end{document}

\end{document}